% SPDX-License-Identifier: CC-BY-4.0
% Copyright (C) 2026 Fabian Januszewski, Abhijit (the Langlandscape team)
% Section 7: main results with their support status (claims table). Statuses are those of the chapters
% (mathematics/chapters/, via final-report/generated/results-table.tex) and of the Lean registry at the final gate; never upgraded.

\paragraph{Status vocabulary.} Each claim below carries one or more of the following statuses, exactly as in the chapters
(where a statement has several clauses, each clause carries its own status).
\begin{itemize}
\item \emph{Proved} (\stproved): the proof is written out in \mainpdf{} \cite{LLmain}; \S\ref{ss:s00-outline} gives its outline.
\item \emph{Proved modulo cited} (\stpmc): proved in \mainpdf, using the cited results named in the statement as black boxes.
\item \emph{Lean-verified}: a Lean~4 declaration with this content compiles at the final gate, its \texttt{\#print axioms} list is a
  subset of \texttt{propext}, \pathtt{Classical.choice}, \texttt{Quot.sound}, it contains no \texttt{sorry}, and its statement
  matches the claim: Table~L1 (\S\ref{ss:lean-tables}) records no difference. We use this word only where the chapter says
  \emph{formalized} without qualifier. Where audited declarations cover only some parts or a special case, we say
  \emph{formalized in part} and Table~L1 records the difference.
\item \emph{Lean-conditional on} \texttt{Hyp.A}, \dots: as Lean-verified, but the declaration takes the registered hypotheses
  \texttt{Hyp.A}, \dots{} as explicit arguments (\S\ref{sec:lean}); it is a theorem \emph{about} those hypotheses, not a proof of them.
\item \emph{Computed}: checked by machine computation over the stated range (\S\ref{sec:comp}); a computation is never part of a proof.
\item \emph{Cited} (\stcited): taken from the literature with a precise pointer; not reproved.
\item \emph{Sketched} (\stsketch), \emph{expected} (\stexp), \emph{conjectural} (\stconj), \emph{open} (\stopen), \emph{to prove} (\sttoprove): no proof is claimed.
\end{itemize}

\paragraph{The headline.} For $D=q$, a $q$-new primitive ordinary Hida branch $F$ with $\lambda_F:\T\to I_F$, $p\ge5$, $\m$
non-Eisenstein with $\bar\rho_\m$ absolutely irreducible on $G_{\Q(\sqrt{p^*})}$ and $p$-distinguished, under the standing
hypotheses (H0), (H1), (H2), (P) and either (H0$^-$) ($q\not\equiv-1\bmod p$ or $\bar\rho_\m$ ramified at $q$) or (Lev)
\mainref{Hyp.~20.1}: the defect $(h_F)=\eta_F(\eta^B_F)^{-1}$ of the congruence ideals of $F$ on the $\GL_2$ side and on the
definite side equals the old/new congruence ideal $\cc_q(F)$, and $\cc_q(F)$ equals the Tamagawa ideal $\tq_q(F)$; at weight-$k$
points $f$ this reads $\eta_f(N)=\cc_q(f)\,\eta^B_f$ with $\cc_q(f)=(\varpi^{t_f(q)})$ \mainref{Thm~21.5, Cor.~21.2, Cor.~21.6}.
Status: \textbf{proved modulo cited} results (Hida's control theorem, Hsieh's Gorenstein theorem, Kim--Ota \cite[Thm~1.1]{KO} and,
for (F1$'$), our sign-fixed reading R4 of \cite[Thm~3.14]{KO}); \textbf{open} for $q\equiv-1\bmod p$ outside (Lev) (i.e.\
$N(\bar\rho)>1$, or non-squarefree $M$); in Lean, \pathtt{Defect.defect\_formula} is \textbf{Lean-conditional} on the registered
hypotheses, including the Galois half \pathtt{Hyp.FreeRankTwo}, and is a statement about an abstract \pathtt{Defect.Setting}
(Hida's Hecke algebra and the Galois representation are not constructed in Lean); its commutative-algebra half (Thm~E) is
formalized as abstract commutative algebra, with the qualifiers of Table~L1 (Lemmas S and G and the cyclic dual lemma are Lean-verified). \textbf{Computed}: 7 triples $(q,k,p)$ at
weight $k$ and 6 of them at weight $k+p-1$ (T5: 4 in stage~06, 2 more by HACS Lane~C), and 225/225 in-scope rows of the regression C2 (\S\ref{sec:comp}).

\paragraph{The claims table.} One row per main claim. ``main.pdf'' gives the number in \mainpdf. In the Lean column, \emph{Lean-verified} means that Table~L1
(\S\ref{ss:lean-tables}) records the statement as in the chapter, with no difference; \emph{formalized in part} means that the
named parts are covered by audited declarations with standard axioms, with the differences (special cases, abstract
settings, extra arguments) recorded in Table~L1; ``cond.'' means Lean-conditional on the hypotheses named; ``hypothesis''
means that the Lean object is the registered hypothesis, i.e.\ the statement is asserted, not proved; ``--'' means not formalized. The last column says what is known and what is ours: \textbf{K} known,
cited from the literature (not reproved); \textbf{S} standard or known, written out in our conventions (no novelty claimed);
\textbf{F} the underlying result or idea is known (reference in the statement), the definite, combinatorial or integral form
given here is ours (we did not find it in this form in the literature we checked); \textbf{N} ours: we did not find this in the
literature we checked; \textbf{C} a hypothesis, conjecture or open item formulated by us. These attributions are our reading of
the literature cited in \mainpdf, not the result of a systematic search; the per-stage account is in the section ``What is new
and what is known''.
Registered hypotheses are listed in Appendix~\ref{app:hyp}; Lean declaration names are in Appendix~\ref{app:decls}.

{\small\setlength{\tabcolsep}{3pt}
\begin{longtable}{@{}>{\raggedright}p{0.25\textwidth} >{\raggedright}p{0.10\textwidth} >{\raggedright}p{0.24\textwidth} >{\raggedright}p{0.26\textwidth} >{\raggedright\arraybackslash}p{0.08\textwidth}@{}}
\caption{Main claims and their support status (statuses as in the chapters and in the Lean registry at the final gate; never upgraded).}\label{tab:claims}\\
\toprule claim & main.pdf & status on paper & Lean; computation & known/ ours\\\midrule\endfirsthead
\toprule claim & main.pdf & status on paper & Lean; computation & known/ ours\\\midrule\endhead
\bottomrule\endfoot
\multicolumn{5}{@{}l}{\textbf{Stage 01: quaternionic foundations}}\\
Eichler mass formula, maximal order (a) & Thm 2.1(a) & cited \cite[Thm~25.1.1]{Vo} & is the hypothesis \pathtt{GL2TF.Hyp.EichlerMass} (cited) & K\\
mass level change (b); indices (c); class number of the tower $|S_s|=\frac{\varphi(D)}{24}\psi(Mp^s)\varphi(Mp^s)$ (d) & Thm 2.1(b)--(d) & (b) proved; (c) proved; (d) proved modulo (a) & (b) ($w$-weighted level change) and (c) first half formalized in part; (d) cond.\ on \pathtt{Hyp.EichlerMass}; computed $|S_s|=2,50,1250$ ($D=3$, $p=5$, $s\le3$) & S\\
norm surjectivity $\Nrd(\cO_\ell^\times)=\Z_\ell^\times$, $\ell\mid D$ & Lemma 1.3 & proved & formalized in part (split primes and assembly); its ramified part \pathtt{Hyp.RamNrdSurj} is \emph{discharged} (\pathtt{Hyp.ramNrdSurj\_holds}, which proves a stronger local statement) & S\\
Brandt pairing, adjoints, Hecke relations, $N_{U_p}$ normalisation & Prop.~3.3, 4.4 & proved / known & formalized in part (Table L1) & S\\
classical JL, Brandt form & Thm 6.1 & cited \cite[Thm~16.1]{JL} & schematic hypothesis \pathtt{Hyp.JLClassical}; computed: 66/66 cases $D\le13$, $k\le12$ & K\\
$p$-adic JL (Chenevier) & Thm 6.2 & cited \cite{Ch} & -- & K\\
$\Lambda$-adic Hecke JL & Thm 6.3 & cited \cite{LV,Hs} & schematic hypothesis \pathtt{Hyp.LambdaJL} & K\\
\multicolumn{5}{@{}l}{\textbf{Stage 02: tree quotients}}\\
strong approximation for $B^1$ at $\{\infty,q\}$ & Thm 7.7 & cited \cite[Main Thm~28.5.3]{Vo}; the appendix proof (Thm 7.34) is proved modulo cited (Hasse--Minkowski) & hypothesis \texttt{Hyp.SA} & K\\
strong approximation at one level $\iff$ the Brandt graph $\BG(U,q)$ is connected (a finite check) & Prop.~7.9 & proved (connectedness of the tree, cited, in one direction) & formalized in part: the tree-free parts; (i)$\Rightarrow$(iii) cond.\ on \texttt{Hyp.BTtree}; computed 240 cases & F\\
$\GU\backslash\cT_q\cong\BG(U,q)$; pull/push maps & Prop.~7.10 & (a) proved given connectivity, proved modulo cited via Thm 7.7 and Hasse--Schilling; (b), (c) proved & (a) formalized given connectivity, or cond.\ on \texttt{Hyp.SA}, \texttt{Hyp.HS} (audited since the submission bundle, Table L1); (b) formalized; (c) formalized only for trivial coefficients and the pull/push maps (not the traces) & S\\
harmonic cochains, $H_1(\GU,V)$ & Prop.~7.18, 7.19 & kernel description proved given connectivity; isomorphism with $H_1(\GU,V)$ proved modulo cited (Serre) & formalized in part ($V=R$, all $w_x$ units); the isomorphism with $H_1$ is the hypothesis \pathtt{Hyp.SerreCoh}; \pathtt{Hyp.SerreFree} & S\\
\multicolumn{5}{@{}l}{\textbf{Stage 03: weighted graphs, Ribet, Ihara}}\\
$\Phicomb$ of two-vertex graphs & Prop.~8.3 & proved & formalized in part ($E\ne\emptyset$, edges oriented $v_0\to v_1$) & F\\
Mazur check: $\Phicomb(\GDR_q(1))\cong\Z/\mathrm{num}\frac{q-1}{12}$ & Prop.~8.8 & proved modulo cited & cond.\ on \pathtt{Hyp.EllipticCount}, \pathtt{Hyp.EichlerMass}; computed $q\in\{11,23,37,101,131\}$ and (HACS) every prime $q\le600$ at $\ell\ge5$ & F\\
$\Phicomb$ is Eisenstein & Prop.~9.5 & proved (with cited inputs) & abstract core formalized; instance at $\hat\cO^\times$ cond.\ on \pathtt{Hyp.BrandtStabFree}, which is discharged (\BSFdischargeName) & F\\
definite level-raising operator & Prop.~9.7 & proved & abstract core formalized in part (its symmetry hypothesis discharged at the quaternionic instances) & F\\
definite Ribet sequence as a vertex collapse & Thm 10.2 & given connectivity of $\BG(U,q)$ (condition (iii)): (a) and the first two sentences of (c) proved; (b) (coefficients) and the last sentence of (c) proved modulo cited; without assuming (iii): proved modulo cited (Thm 7.7, Hasse--Schilling, connectedness of the tree) & (a) and the rank in (c) formalized given connectivity (\pathtt{Tree.BGConnected}), or cond.\ on \texttt{Hyp.SA}, \texttt{Hyp.HS}, \texttt{Hyp.BTtree}; (b) -- & F\\
definite Ihara lemma (i), (ii) & Cor.~10.5 & (i), (ii) proved given connectivity (condition (iii)); (ii$'$) proved modulo cited (given (iii); the tree property is cited) & (i), (ii) formalized over any ring, given connectivity (or cond.\ on \texttt{Hyp.SA}, \texttt{Hyp.HS}, \texttt{Hyp.BTtree}); (ii$'$) -- & F\\
Ihara (iii): $H_0$ and its Hecke action & Thm 11.3 & proved modulo cited & clause (a) is the hypothesis \texttt{Hyp.IharaH0}; (b) -- & F\\
Eisenstein support & Cor.~11.5 & proved modulo cited & -- & F\\
amalgam, Mayer--Vietoris, $\det\Psi^*\Psi=(q+1)^2-q^{2-k}T_q^2$ & Prop.~12.2, 12.3, 12.6 & proved modulo cited & --; computed $\det\Psi^*\Psi$ in 121 cases (never $0$) & S\\
Ihara lemma for $\GL_2$, Betti form & Conj.~12.10 & expected (cited for parts) & -- & K\\
\multicolumn{5}{@{}l}{\textbf{Stage 04: ordinary and $\Lambda$-adic}}\\
freeness at finite level & Prop.~14.2 & proved & formalized in part ($\Zp$, $U=\Uone M$, $p\ge5$) & S\\
$\Lambda$-freeness and control of $\Mord(U)$ & Thm 14.4 & (a) proved; (b) cited (integrally if $p\nmid\varphi(M)$, otherwise sketched), general level expected & (a) finite-level part formalized ($\Zp$, $\Uone M$, $p\ge5$); the $\Lambda$-adic half is the open hypothesis \pathtt{Hyp.MordLambdaFree} (stated only; no Lean theorem uses it); (b) schematic \pathtt{Hyp.HidaControl} & F\\
$\Lambda$-adic definite Ribet sequence & Thm 15.1 & proved, given connectivity of $\BG(UU_s,q)$ for all $s\ge1$ (a hypothesis of the theorem) & finite levels $s\ge1$ formalized ($\Qp$, $\Uone M$, $p\ge5$); inverse limit formalized given connectivity and the surjectivity of the $A$-transitions as explicit arguments (the latter is the conclusion of the open \pathtt{Hyp.InvLimExact}) & F\\
Hsieh: Gorenstein and freeness; $\eta^B_F$ as a self-pairing & Thm 16.6, 16.7 & (i)--(iii) cited \cite[Thm~4.5]{Hs}, (iv) proved from (ii); 16.7 cited & hypothesis \pathtt{Hyp.HsiehGor} & K\\
$\eta$-principle & Prop.~17.1 & proved & \textbf{Lean-verified} (\pathtt{CommAlg.EtaPrinciple}) & S\\
$e$ kills $H_0$ with $\Vk$-coefficients & Prop.~18.1 & \textbf{to prove} (transfer step not written) & -- & C\\
\multicolumn{5}{@{}l}{\textbf{Stage 05: the comparison hypothesis}}\\
$\HypX(M;q)$: Hecke-equivariant isometries $\Xq(J_s)\cong\Z[S_s]^0$ compatible in $s$ and with Galois & Hyp.~18.15 & a \textbf{hypothesis}, open at $\Gamma_1(Mp^s)$-level with $U_p$; never assumed known & hypotheses \texttt{Hyp.X.char}, \texttt{Hyp.X.iso}, \texttt{Hyp.X.gal} & C\\
(G1)--(G4) at $\Gamma_0$-type level (Ribet, Deligne--Rapoport, Grothendieck) & Thm 18.30 & cited \cite{Ri90,DR,SGA7} & -- & K\\
sign $\varepsilon=+1$, exponent $+1$: $U_q\mapsto+\Pi_q$ & Prop.~18.6 & proved modulo cited (given (X1) and the model choice FD1 at $\Gamma_1$-level); the hybrid dictionary excluded only modulo cited results at suitable levels & --; computed $a_q=+\lambda(\Pi_q)$ on 144/144 eigenlines & F\\
$\Lambda$-adic version of $\HypX$ & Prop.~18.20 & parts proved / proved modulo cited; (c2) conjectural, open & \pathtt{xlambda\_finite} cond.\ on \texttt{Hyp.X.char}, for an abstract \texttt{CharGroupData}, existence only & F\\
$\HypXiso$ holds non-canonically & Prop.~18.24 & proved modulo cited (parts sketched) & algebraic half formalized (\pathtt{xiso\_of\_gorenstein}) & F\\
$q$-new sub- versus quotient lattice & Prop.~18.25 & proved & \textbf{Lean-verified} (\pathtt{CommAlg.subquot\_kernel}, \pathtt{subquot\_coker}), with the splitting hypothesis in the equivalent form $(\fa+\fb)\T_K=\T_K$ (the CRT equivalence with $\T_K=(\T/\fa)_K\times(\T/\fb)_K$ is not formalized) & F\\
no tree-only chain & Prop.~18.28 & proved / proved modulo cited & -- & N\\
relation to Chenevier (O7) & Rem.~18.41 & sketched on the ordinary locus; open beyond & -- & F\\
simplicial realisation O2 & Conj.~18.40 & conjectural, open & -- & C\\
\multicolumn{5}{@{}l}{\textbf{Stage 06: the integral Jacquet--Langlands defect}}\\
Lemma S, Lemma G, cyclic dual lemma & Lemma 19.2--19.4 & proved & \textbf{Lean-verified} (full statements) & S\\
Lemma D (new/old data) & Lemma 19.5 & proved & formalized with qualifiers (splitting hypothesis field-free; closing base-change sentence not formalized) & F\\
Thm E, the combinatorial half: $\eta=\cc\cdot\eta'$ for new/old data & Thm 19.8 & proved & (i)--(iv) formalized ((iv) for $S'$ $A$-free) as abstract commutative algebra (\texttt{NewOldData}); not formalized: localisation clause of (iii), specialisation to Hecke rings & F\\
Kim--Ota & Thm 21.1 & cited \cite[Thm~1.1]{KO} & hypothesis \texttt{Hyp.KimOta} & K\\
fixed-weight level lowering mod $\varpi^t$ & Cor.~21.2 & proved modulo cited & \texttt{cor\_fixed} cond.\ on registered hypotheses & F\\
\textbf{$\Lambda$-adic defect formula} $(h_F)=\cc_q(F)=\tq_q(F)$ & \textbf{Thm 21.5} & \textbf{proved modulo cited} under Hyp.~20.1; \textbf{open} for $q\equiv-1\bmod p$ outside (Lev) & \pathtt{Defect.defect\_formula} cond.\ on the registered hypotheses incl.\ \pathtt{Hyp.FreeRankTwo} (abstract \pathtt{Defect.Setting}); computed T5, C2 & N\\
all weights of the branch & Cor.~21.6 & proved modulo cited & \pathtt{cor\_allweights} cond.\ (the newform condition $\cc\ne0$ is not a Lean argument) & N\\
Galois comparison $\tq_q\subseteq\cc_q$ & Prop.~22.1 & proved modulo cited & \pathtt{Defect.galois}, \pathtt{freeRankTwo\_iff\_le}: given \pathtt{Hyp.ResidualBurnside}, \pathtt{Hyp.FreeRankTwo} $\iff$ the inclusion; the Galois half is not formalized independently & F\\
the inclusion without (H0$^-$), $q\equiv-1\bmod p$ & Prop.~23.2 & (F0), (F1) proved modulo cited; (F1$'$) proved modulo cited in our reading R4 of \cite[Thm~3.14]{KO}; $N(\bar\rho)>1$, non-squarefree $M$ open & hypothesis \pathtt{Hyp.KimOtaRT}; computed (H3): numerical evidence in the open region; (F1) untested & F\\
Betti incarnation of $\cc_q$ & Prop.~24.1 & (a) proved; (b), (c) proved modulo cited & --; computed T4 ($q=3$, $p=17$) & F\\
several ramified primes $(h_F)=\prod_i\tq_{q_i}$ & Conj.~25.2 & conjectural & -- & C\\
\multicolumn{5}{@{}l}{\textbf{Literature}}\\
sign correction $\JL(U_q)=+U_q$ for $q\mid N^-$ in \cite[\S4.5]{Hs} & \S\ref{ss:hsieh} & proved modulo cited for trivial central character at $q$; sketched otherwise & --; computed 144/144 eigenlines & N\\
\end{longtable}}

\noindent The complete list of the 88 theorems, propositions, corollaries and conjectures of \mainpdf{} with all status markers
is reproduced (generated from the chapter sources) in Table~\ref{tab:allresults}.

{\small
% SPDX-License-Identifier: CC-BY-4.0
% Copyright (C) 2026 Fabian Januszewski, Abhijit (the Langlandscape team)
% GENERATED by tools/extract_statements.py from mathematics/chapters/ -- do not edit by hand.

\begin{longtable}{@{}>{\raggedright}p{0.13\textwidth} >{\raggedright}p{0.25\textwidth} >{\raggedright}p{0.30\textwidth} >{\raggedright\arraybackslash}p{0.25\textwidth}@{}}
\caption{Every theorem, proposition, corollary and conjecture of \texttt{mathematics/main.pdf} (number, page), with the status markers that occur inside its statement and, marked ``after the statement'', those of a status paragraph printed directly after it in the chapter (generated; \texttt{\textbackslash lean} refs = number of Lean cross-references in the statement). Mixed markers refer to different clauses; the statement in main.pdf is authoritative. A ``--'' means that no marker was found in either place, not that the result is proved. Rows are in the numerical order of main.pdf within each stage.}\label{tab:allresults}\\
\toprule no.\ (p.) & label & title & status markers\\\midrule\endfirsthead
\toprule no.\ (p.) & label & title & status markers\\\midrule\endhead
\multicolumn{4}{@{}l}{\textbf{Stage 01}}\\
Prop. 1.4 (p.~28) & \texttt{\footnotesize prop:s01-eichler} & Eichler orders & \stknown{}\\
Prop. 1.6 (p.~29) & \texttt{\footnotesize prop:s01-finite} & finiteness & \stknown{}, \stformal{}; 2 \texttt{\textbackslash lean} refs\\
Theo. 2.1 (p.~31) & \texttt{\footnotesize thm:s01-mass} & Eichler mass formula; \texttt{if:s01-mass} & \stcited{}, \stformal{}, \stproved{}, \stpmc{}; 1 \texttt{\textbackslash lean} refs\\
Prop. 3.3 (p.~33) & \texttt{\footnotesize prop:s01-adjoint} & adjoints; equivariance of $\nu$ & \stcited{}, \stproved{}, \stformal{}; 8 \texttt{\textbackslash lean} refs\\
Prop. 4.4 (p.~39) & \texttt{\footnotesize prop:s01-relations} & relations & \stproved{}, \stformal{}, \stknown{}; 4 \texttt{\textbackslash lean} refs\\
Theo. 6.1 (p.~48) & \texttt{\footnotesize thm:s01-JLclassical} & classical JL, Brandt form; \pathtt{if:s01-targets}; seed Thm.~3.1 & \stcited{}; after the statement: \stformal{}\\
Theo. 6.2 (p.~49) & \texttt{\footnotesize thm:s01-chenevier} & $p$-adic JL; Chenevier & \stcited{}\\
Theo. 6.3 (p.~49) & \texttt{\footnotesize thm:s01-lambdaJL} & $\Lambda$-adic Hecke JL & \stcited{}, \textsf{\small[provisional note]}; 1 \texttt{\textbackslash lean} refs; after the statement: \stformal{}\\
\multicolumn{4}{@{}l}{\textbf{Stage 02}}\\
Theo. 7.7 (p.~54) & \texttt{\footnotesize thm:s02-SA} & strong approximation for $B^1$ at $S=\{\infty,q\}$ & after the statement: \stcited{}, \stproved{}\\
Prop. 7.9 (p.~55) & \texttt{\footnotesize prop:s02-BG} & SA for one level as a finite check & after the statement: \stproved{}, \stcited{}, \stformal{}\\
Prop. 7.10 (p.~57) & \texttt{\footnotesize prop:s02-quot} & quotient graph; seed Prop.~quot & \stcited{}; after the statement: \stproved{}, \stpmc{}\\
Prop. 7.14 (p.~61) & \texttt{\footnotesize prop:s02-Mhomprops} & properties of Def.~7.13 & after the statement: \stproved{}, \stformal{}\\
Prop. 7.16 (p.~64) & \texttt{\footnotesize prop:s02-duality} & the two pairings & \stformal{}; 2 \texttt{\textbackslash lean} refs; after the statement: \stproved{}\\
Prop. 7.18 (p.~65) & \texttt{\footnotesize prop:s02-coh} & seed Prop.~coh & after the statement: \stcited{}, \stpmc{}, \stproved{}, \stformal{}\\
Prop. 7.19 (p.~66) & \texttt{\footnotesize prop:s02-harm} & harmonic cochains; seed Prop.~harm & after the statement: \stproved{}, \stpmc{}, \stformal{}\\
Prop. 7.20 (p.~67) & \texttt{\footnotesize prop:s02-free} & the free case & after the statement: \stpmc{}, \stproved{}, \stformal{}\\
Prop. 7.22 (p.~68) & \texttt{\footnotesize prop:s02-indefprops} & properties of the tree model & after the statement: \stpmc{}, \stproved{}, \stformal{}\\
Prop. 7.31 (p.~73) & \texttt{\footnotesize prop:s02-sa-step-d} & tex L02.SA(d) & \stpmc{}, \stcited{}\\
Theo. 7.34 (p.~74) & \texttt{\footnotesize thm:s02-sa-main} & Thm.~7.7, proof & \stpmc{}\\
\multicolumn{4}{@{}l}{\textbf{Stage 03}}\\
Prop. 8.3 (p.~76) & \texttt{\footnotesize prop:s03-phi} & two-vertex graphs; seed Prop.~phi & \stproved{}, \stformal{}; 4 \texttt{\textbackslash lean} refs\\
Prop. 8.8 (p.~79) & \texttt{\footnotesize prop:s03-mazur} & Mazur check; seed Example mazur & \stpmc{}, \stformal{}; 8 \texttt{\textbackslash lean} refs\\
Prop. 9.5 (p.~83) & \texttt{\footnotesize prop:s03-phieis} & $\Phicomb$ is Eisenstein & \stproved{}, \stformal{}, \stcited{}; 4 \texttt{\textbackslash lean} refs\\
Prop. 9.7 (p.~85) & \texttt{\footnotesize prop:s03-LRdef} & definite level-raising operator & \stproved{}, \stformal{}, \stcited{}; 11 \texttt{\textbackslash lean} refs\\
Theo. 10.2 (p.~89) & \texttt{\footnotesize thm:s03-ribet} & definite Ribet sequence as a vertex collapse; seed Thm ribet & \stproved{}, \stpmc{}, \stcited{}, \stformal{}, \leanplanned{}; 14 \texttt{\textbackslash lean} refs\\
Coro. 10.5 (p.~90) & \texttt{\footnotesize cor:s03-ihara} & definite Ihara lemma, (i)--(ii); seed Cor.~ihara & \stproved{}, \stpmc{}, \stformal{}; 4 \texttt{\textbackslash lean} refs\\
Theo. 11.3 (p.~95) & \texttt{\footnotesize thm:s03-iharaH0} & Ihara (iii): $H_0$ and its Hecke action & \stpmc{}, \sttoprove{}, \stsketch{}, \stproved{}, \stformal{}, \leanplanned{}; 5 \texttt{\textbackslash lean} refs\\
Coro. 11.5 (p.~98) & \texttt{\footnotesize cor:s03-eis} & Eisenstein support & \stpmc{}, \stcited{}, \leanplanned{}\\
Prop. 12.2 (p.~101) & \texttt{\footnotesize prop:s03-amalgam} & amalgam & \stpmc{}, \stcited{}\\
Prop. 12.3 (p.~101) & \texttt{\footnotesize prop:s03-MV} & Mayer--Vietoris; seed \S6.4 & \stpmc{}, \stcited{}\\
Prop. 12.6 (p.~104) & \texttt{\footnotesize prop:s03-PsiPsi} & $\Psi^*\Psi$ & \stpmc{}, \sttoprove{}, \stcomp{}\\
Conj. 12.10 (p.~107) & \texttt{\footnotesize conj:s03-iharaGL2} & Ihara lemma for $\GL_2$, Betti form & \stexp{}, \stcited{}\\
\multicolumn{4}{@{}l}{\textbf{Stage 04}}\\
Prop. 14.2 (p.~112) & \texttt{\footnotesize prop:s04-freeA} & freeness at finite level; seed Prop.~free (a) & \stproved{}\\
Theo. 14.4 (p.~114) & \texttt{\footnotesize thm:s04-control} & $\Lambda$-freeness and control; seed Prop.~free (b); \texttt{if:s04-Mord} & \stexp{}, \stsketch{}, \stproved{}, \stpmc{}, \stcited{}, \stformal{}, \textsf{\small[provisional note]}; 7 \texttt{\textbackslash lean} refs\\
Theo. 15.1 (p.~116) & \texttt{\footnotesize thm:s04-lambdaribet} & $\Lambda$-adic definite Ribet sequence; seed Thm.~lambdaribet; \pathtt{if:s04-lambdaribet} & \stproved{}, \stformal{}, \textsf{\small[provisional note]}; 5 \texttt{\textbackslash lean} refs\\
Theo. 16.6 (p.~122) & \texttt{\footnotesize thm:s04-Gor} & Gorenstein and freeness; \texttt{if:s04-Mord}, second half & \stexp{}; after the statement: \stcited{}, \stproved{}\\
Theo. 16.7 (p.~123) & \texttt{\footnotesize thm:s04-etaB} & $\eta^B_F$ as a self-pairing; \texttt{if:s04-etaB} & after the statement: \stcited{}\\
Prop. 17.1 (p.~124) & \texttt{\footnotesize prop:s04-etaprinciple} & $\eta$-principle; \pathtt{thm:s04-etaB}, \texttt{lem:s06-S} & \stproved{}, \stformal{}; 4 \texttt{\textbackslash lean} refs\\
Prop. 18.1 (p.~125) & \texttt{\footnotesize prop:s04-ekillsVk} & $e$ kills $H_0$ with $\Vk$-coefficients & after the statement: \sttoprove{}\\
\multicolumn{4}{@{}l}{\textbf{Stage 05}}\\
Prop. 18.6 (p.~128) & \texttt{\footnotesize prop:s05-sign} & $\varepsilon=+1$; exponent $+1$ & \stpmc{}, \stproved{}, \stknown{}, \sttoprove{}, \sthyp{}, \textsf{\small[provisional note]}\\
Prop. 18.13 (p.~137) & \texttt{\footnotesize prop:s05-signe} & Prop.~sign (e): higher weight & \stpmc{}, \stexp{}\\
Prop. 18.18 (p.~141) & \texttt{\footnotesize prop:s05-AL} & Atkin--Lehner variant & \stpmc{}, \stcited{}, \stexp{}, \stsketch{}, \textsf{\small[provisional note]}\\
Prop. 18.20 (p.~144) & \texttt{\footnotesize prop:s05-Xlambda} & $\Lam$-adic version; seed Prop.~Xlambda & \stproved{}, \stpmc{}, \stopen{}, \stsketch{}, \stconj{}, \textsf{\small[provisional note]}; after the statement: \stformal{}, \leanplanned{}\\
Prop. 18.24 (p.~150) & \texttt{\footnotesize prop:s05-Xiso} & $\HypXiso$ holds non-canonically & \stpmc{}, \stformal{}, \stsketch{}; 1 \texttt{\textbackslash lean} refs\\
Prop. 18.25 (p.~151) & \texttt{\footnotesize prop:s05-subquot} & $q$-new sub- versus quotient lattice; seed Prop.~subquot; \pathtt{if:s05-subquot} & \stproved{}, \stformal{}; 2 \texttt{\textbackslash lean} refs\\
Prop. 18.28 (p.~152) & \texttt{\footnotesize prop:s05-nochain} & no tree-only chain; corrected form of seed Prop.~noChain & \stproved{}, \stpmc{}\\
Theo. 18.30 (p.~153) & \texttt{\footnotesize thm:s05-G} & (G1)--(G4) at $\Gamma_0$-type level & \stcited{}; after the statement: \stcited{}\\
Conj. 18.40 (p.~159) & \texttt{\footnotesize conj:s05-O2} & O2, a simplicial realisation; open & \stconj{}, \stopen{}\\
\multicolumn{4}{@{}l}{\textbf{Stage 06}}\\
Theo. 19.8 (p.~163) & \texttt{\footnotesize thm:s06-E} & the combinatorial half; seed Thm~E, hypotheses explicit & \stproved{}, \stformalized{}; 7 \texttt{\textbackslash lean} refs\\
Theo. 21.1 (p.~170) & \texttt{\footnotesize thm:s06-KO} & Kim--Ota & \stcited{}\\
Coro. 21.2 (p.~170) & \texttt{\footnotesize cor:s06-fixed} & fixed-weight level lowering mod $\varpi^t$; replaces seed step (4) & \stpmc{}\\
Theo. 21.5 (p.~172) & \texttt{\footnotesize thm:s06-defect} & $\Lam$-adic defect formula; seed Thm~defect, corrected & \stpmc{}, \stopen{}\\
Coro. 21.6 (p.~173) & \texttt{\footnotesize cor:s06-allweights} & all weights of the branch; seed Cor.~allweights & \stpmc{}\\
Prop. 22.1 (p.~174) & \texttt{\footnotesize prop:s06-galois} & Galois comparison; seed Prop.~galois & \stpmc{}; after the statement: \stcited{}, \stknown{}\\
Prop. 23.2 (p.~176) & \texttt{\footnotesize prop:s06-levM} & the inclusion without (H0$^-$); the seed's step (4) & \stpmc{}, \stopen{}\\
Prop. 24.1 (p.~184) & \texttt{\footnotesize prop:s06-betti} & Betti incarnation of $\cc_q$ & \stproved{}, \stpmc{}, \textsf{\small[provisional note]}\\
Conj. 25.2 (p.~187) & \texttt{\footnotesize conj:s06-several} & several ramified primes; seed \S10, last paragraph & \stconj{}, \stsketch{}, \stopen{}\\
\bottomrule\end{longtable}
}
